\documentclass[11pt]{article}
\usepackage[margin=1in]{geometry}
\usepackage{listings,xcolor,hyperref,booktabs,longtable,graphicx}
\lstset{basicstyle=\ttfamily\small,breaklines=true,frame=single,numbers=left,numberstyle=\tiny\color{gray}}
\title{Generative UI Architecture:\\SAP Analytics Cloud AI Widget System}
\author{SAP OSS Technical Documentation}
\date{March 2026}

\hypersetup{
  colorlinks=true,
  linkcolor=blue!70!black,
  urlcolor=blue!70!black,
  citecolor=blue!70!black,
}

\begin{document}

\maketitle

\begin{abstract}
This document provides a comprehensive technical analysis of the Generative UI
subsystem within the SAP OSS codebase. The system enables an AI-driven widget
experience inside SAP Analytics Cloud (SAC) by combining a 13-type widget schema,
an AG-UI protocol bridge over Server-Sent Events, a tool dispatch service with
human-in-the-loop approval, and a Python backend with circuit-breaker resilience.
All source evidence references exact file paths and line numbers from the
repository at \texttt{src/generativeUI/}.
\end{abstract}

\tableofcontents
\newpage

% =============================================================================
\section{Widget Schema Architecture}
\label{sec:schema}
% =============================================================================

The widget schema is defined in
\texttt{libs/sac-ai-widget/types/sac-widget-schema.ts} and establishes a
13-type taxonomy organized into five categories.

\subsection{Two-Level Type Taxonomy}

The type system uses a union-of-unions pattern. Four subordinate type aliases
feed into the top-level \texttt{SacWidgetType} union (lines~11--31):

\begin{lstlisting}[language={[Sharp]C},caption={Widget type taxonomy (sac-widget-schema.ts:11--31)}]
export type SacWidgetTypeCore   = 'chart' | 'table' | 'kpi';
export type SacWidgetTypeFilter = 'filter-dropdown'
  | 'filter-checkbox' | 'filter-date-range';
export type SacWidgetTypeSlider = 'slider' | 'range-slider';
export type SacWidgetTypeText   = 'text-block' | 'heading' | 'divider';
export type SacWidgetTypeLayout = 'grid-container' | 'flex-container';

export type SacWidgetType =
  | SacWidgetTypeCore
  | SacWidgetTypeFilter
  | SacWidgetTypeSlider
  | SacWidgetTypeText
  | SacWidgetTypeLayout;
\end{lstlisting}

\noindent The categories and their roles are:

\begin{table}[h]
\centering
\begin{tabular}{lll}
\toprule
\textbf{Category} & \textbf{Types} & \textbf{Role} \\
\midrule
Core         & chart, table, kpi                          & Data visualization \\
Filter       & filter-dropdown, filter-checkbox, filter-date-range & Interactive filtering \\
Slider       & slider, range-slider                       & Numeric range input \\
Text         & text-block, heading, divider               & Content display \\
Layout       & grid-container, flex-container              & Container/composition \\
\bottomrule
\end{tabular}
\caption{Widget type taxonomy.}
\label{tab:types}
\end{table}

\subsection{Runtime Validation}

The \texttt{VALID\_WIDGET\_TYPES} \texttt{Set<SacWidgetType>} (lines~34--40)
provides O(1) membership testing. The \texttt{validateWidgetType()} type guard
function (lines~46--48) combines a \texttt{typeof} check with set membership:

\begin{lstlisting}[language={[Sharp]C},caption={Runtime type guard (sac-widget-schema.ts:34--48)}]
export const VALID_WIDGET_TYPES = new Set<SacWidgetType>([
  'chart', 'table', 'kpi',
  'filter-dropdown', 'filter-checkbox', 'filter-date-range',
  'slider', 'range-slider',
  'text-block', 'heading', 'divider',
  'grid-container', 'flex-container',
]);

export const MAX_CHILDREN_DEPTH = 8;

export function validateWidgetType(
    value: unknown): value is SacWidgetType {
  return typeof value === 'string'
    && VALID_WIDGET_TYPES.has(value as SacWidgetType);
}
\end{lstlisting}

\noindent The \texttt{MAX\_CHILDREN\_DEPTH = 8} constant (line~43) bounds
recursion when normalizing nested container schemas.

\subsection{Type Guards}

Five boolean type guard functions (lines~189--207) classify a
\texttt{SacWidgetSchema} by its \texttt{widgetType} field:

\begin{itemize}
  \item \texttt{isContainerWidget()} -- returns \texttt{true} for
        \texttt{grid-container} and \texttt{flex-container} (line~190).
  \item \texttt{isFilterWidget()} -- prefix match on \texttt{filter-}
        (line~194).
  \item \texttt{isSliderWidget()} -- matches \texttt{slider} and
        \texttt{range-slider} (line~198).
  \item \texttt{isTextWidget()} -- matches \texttt{text-block},
        \texttt{heading}, and \texttt{divider} (line~202).
  \item \texttt{isDataWidget()} -- matches core types \texttt{chart},
        \texttt{table}, \texttt{kpi} (line~206).
\end{itemize}

\subsection{Schema Interface}

The \texttt{SacWidgetSchema} interface (lines~140--176) unifies all 13 widget
types into a single shape. Data widgets require \texttt{modelId},
\texttt{dimensions}, and \texttt{measures}. Container widgets use the
\texttt{children} array. Filter, slider, text, and layout widgets use their
respective configuration sub-interfaces (\texttt{SacSliderConfig},
\texttt{SacTextConfig}, \texttt{SacLayoutConfig}/\texttt{SacGridConfig}).
Accessibility properties \texttt{ariaLabel} and \texttt{ariaDescription} appear
on every widget type (lines~173--175).

% =============================================================================
\section{AG-UI Protocol Bridge}
\label{sec:agui}
% =============================================================================

The \texttt{SacAgUiService} (file:
\texttt{libs/sac-ai-widget/ag-ui/sac-ag-ui.service.ts}) bridges the Angular
client to the CAP LLM Plugin backend using the AG-UI protocol over Server-Sent
Events.

\subsection{Event Type Enumeration}

Seventeen event types are defined as a string union (lines~27--33):

\begin{lstlisting}[language={[Sharp]C},caption={AG-UI event types (sac-ag-ui.service.ts:27--33)}]
export type AgUiEventType =
  | 'RUN_STARTED' | 'RUN_FINISHED' | 'RUN_ERROR'
  | 'STEP_STARTED' | 'STEP_FINISHED'
  | 'TEXT_MESSAGE_START' | 'TEXT_MESSAGE_CONTENT'
  | 'TEXT_MESSAGE_END'
  | 'TOOL_CALL_START' | 'TOOL_CALL_ARGS'
  | 'TOOL_CALL_END' | 'TOOL_CALL_RESULT'
  | 'STATE_SNAPSHOT' | 'STATE_DELTA'
  | 'MESSAGES_SNAPSHOT'
  | 'CUSTOM' | 'RAW';
\end{lstlisting}

\noindent A \texttt{Set<AgUiEventType>} named \texttt{KNOWN\_EVENT\_TYPES}
(lines~101--119) mirrors this union for runtime validation during SSE frame
parsing.

\subsection{SSE Streaming}

The \texttt{run()} method (lines~148--238) performs a \texttt{POST} to
\texttt{/ag-ui/run} with \texttt{Accept: text/event-stream} and returns an
RxJS \texttt{Observable<AgUiEvent>}. Key implementation details:

\begin{itemize}
  \item An \texttt{AbortController} (line~150) is created per run and tracked
        in the \texttt{activeControllers} set (line~130) for lifecycle cleanup.
  \item The SSE stream is consumed via the Fetch API's
        \texttt{ReadableStream} reader (lines~186--214).
  \item Frame parsing handles \texttt{CRLF} normalization, size limits
        (\texttt{MAX\_SSE\_FRAME\_LENGTH} = 64\,KB, line~121;
        \texttt{MAX\_SSE\_BUFFER\_LENGTH} = 256\,KB, line~122), and the
        \texttt{[DONE]} sentinel (line~323).
  \item Unsubscribing from the Observable calls \texttt{controller.abort()}
        (line~234), which cleanly terminates the SSE connection. Abort errors
        are silently swallowed (line~225).
\end{itemize}

\subsection{W3C Traceparent Correlation}

Both \texttt{run()} and \texttt{dispatchToolResult()} construct a
\texttt{traceparent} header conforming to W3C Trace Context format
(lines~165--171, 256--261):

\begin{lstlisting}[language={[Sharp]C},caption={Traceparent construction (sac-ag-ui.service.ts:165--171)}]
const traceId = this.sessionService.getTraceId();
const spanId = traceId.slice(0, 16);
const headers: Record<string, string> = {
  'Content-Type': 'application/json',
  Accept: 'text/event-stream',
  traceparent: `00-${traceId}-${spanId}-01`,
};
\end{lstlisting}

\subsection{Tool Result Dispatch}

The \texttt{dispatchToolResult()} method (lines~243--288) posts completed
frontend tool results back to \texttt{/ag-ui/tool-result}. It validates the
\texttt{toolCallId} against the \texttt{pendingToolCalls} map (line~244),
prevents duplicate submissions via the \texttt{inFlightToolResults} set
(line~250), and includes the same traceparent correlation header.

\subsection{Lifecycle Cleanup}

The \texttt{ngOnDestroy()} implementation (lines~290--297) aborts all active
controllers and clears the pending and in-flight tracking sets.

% =============================================================================
\section{Tool Dispatch Service}
\label{sec:tools}
% =============================================================================

The \texttt{SacToolDispatchService} (file:
\texttt{libs/sac-ai-widget/chat/sac-tool-dispatch.service.ts}) executes
frontend-only SAC tool calls dispatched by the LLM via AG-UI events.

\subsection{Five Tool Handlers}

The \texttt{execute()} method (lines~113--128) dispatches to five named tools:

\begin{lstlisting}[language={[Sharp]C},caption={Tool dispatch switch (sac-tool-dispatch.service.ts:113--128)}]
async execute(toolName: string,
    args: Record<string, unknown>): Promise<ToolResult> {
  switch (toolName) {
    case 'set_datasource_filter':
      return this.setDatasourceFilter(args);
    case 'set_chart_type':
      return this.setChartType(args);
    case 'run_data_action':
      return this.runDataAction(args);
    case 'get_model_dimensions':
      return this.getModelDimensions(args);
    case 'generate_sac_widget':
      return this.generateSacWidget(args);
    default:
      return { success: false,
               error: `Unknown tool: ${toolName}` };
  }
}
\end{lstlisting}

\begin{description}
  \item[\texttt{set\_datasource\_filter}] (lines~130--159) Applies a dimension
    filter with type normalization (\texttt{SingleValue}, \texttt{MultipleValue},
    \texttt{RangeValue}, \texttt{AllValue}, \texttt{ExcludeValue}) and refreshes
    the active widget target.
  \item[\texttt{set\_chart\_type}] (lines~161--177) Mutates the chart type on
    the active widget and triggers a data refresh.
  \item[\texttt{run\_data\_action}] (lines~179--235) Executes a SAC planning
    data action. Supports foreground and background execution modes. Requires
    human approval via \texttt{getConfirmationReview()}.
  \item[\texttt{get\_model\_dimensions}] (lines~286--309) Returns the current
    binding metadata (model, dimensions, measures, filters, chart type).
  \item[\texttt{generate\_sac\_widget}] (lines~311--361) Creates a complete
    widget from LLM-generated schema. Validates \texttt{widgetType} via
    \texttt{validateWidgetType()} (line~315), enforces \texttt{modelId} for
    data widgets (lines~324--332), and parses all sub-configurations.
\end{description}

\subsection{ToolExecutionReview with Risk Assessment}

The \texttt{ToolExecutionReview} interface (lines~64--76) structures the
human-in-the-loop approval payload:

\begin{lstlisting}[language={[Sharp]C},caption={Review interface (sac-tool-dispatch.service.ts:64--76)}]
export interface ToolExecutionReview {
  toolName: string;
  title: string;
  summary: string;
  confirmationLabel: string;
  riskLevel: ToolReviewRiskLevel; // 'medium' | 'high'
  actionId?: string;
  modelId?: string;
  binding?: WidgetBindingInfo;
  normalizedArgs: Record<string, unknown>;
  affectedScope: string[];
  rollbackPreview: ToolRollbackPreview;
}
\end{lstlisting}

\noindent The \texttt{buildDataActionReview()} method (lines~237--284)
constructs this review for \texttt{run\_data\_action} calls, including planning
version info and lock status from \texttt{SacPlanningModelService}.

\subsection{Rollback Preview}

The \texttt{ToolRollbackPreview} interface (lines~58--62) provides two
strategies: \texttt{revertData} for foreground actions (line~578) and
\texttt{manualReview} for background executions (line~570). Each includes a
human-readable \texttt{label} and contextual \texttt{warnings} about working
version state and lock status.

\subsection{Widget State Target Pattern}

The service uses a \texttt{WidgetStateTarget} interface (lines~50--54) that
the \texttt{SacAiDataWidgetComponent} implements. Registration occurs via
\texttt{registerTarget()} (line~93) and \texttt{unregisterTarget()} (line~97),
enabling tool methods to call \texttt{applySchema()} and \texttt{refreshData()}
on the active widget.

% =============================================================================
\section{Data Widget Component}
\label{sec:datawidget}
% =============================================================================

The \texttt{SacAiDataWidgetComponent} (file:
\texttt{libs/sac-ai-widget/data-widget/sac-ai-data-widget.component.ts}) renders
all 13 widget types and manages live data binding.

\subsection{Component Configuration}

The component is declared as Angular standalone with \texttt{OnPush} change
detection (line~84). It self-references via \texttt{forwardRef()} (line~82) to
support recursive rendering of container children:

\begin{lstlisting}[language={[Sharp]C},caption={Component metadata (sac-ai-data-widget.component.ts:66--84)}]
@Component({
  selector: 'sac-ai-data-widget',
  standalone: true,
  imports: [
    CommonModule, SacChartModule, SacTableModule,
    SacAdvancedModule,
    SacFilterDropdownComponent, SacFilterCheckboxComponent,
    SacSliderComponent,
    SacHeadingComponent, SacTextBlockComponent,
    SacDividerComponent,
    SacFlexContainerComponent, SacGridContainerComponent,
    forwardRef(() => SacAiDataWidgetComponent),
  ],
  changeDetection: ChangeDetectionStrategy.OnPush,
  ...
})
\end{lstlisting}

\subsection{13 Widget Type Renderings}

The template uses \texttt{ngSwitch} on \texttt{schema.widgetType} (line~96) to
select among 13 rendering branches:

\begin{enumerate}
  \item \textbf{chart} (lines~97--107): \texttt{<sac-chart>} with computed
    feeds, chart type, legend position.
  \item \textbf{table} (lines~109--117): \texttt{<sac-table>} with dynamic
    columns and rows.
  \item \textbf{kpi} (lines~119--125): \texttt{<sac-kpi>} with trend
    direction (up/down/neutral).
  \item \textbf{heading} (lines~127--132): \texttt{<sac-heading>} with
    semantic level (h1--h6).
  \item \textbf{text-block} (lines~134--139): \texttt{<sac-text-block>} with
    optional markdown rendering.
  \item \textbf{divider} (lines~141--144): \texttt{<sac-divider>} with ARIA
    separator role.
  \item \textbf{filter-dropdown} (lines~146--157): \texttt{<sac-filter-dropdown>}
    with single/multi select.
  \item \textbf{filter-checkbox} (lines~159--167): \texttt{<sac-filter-checkbox>}
    with checkbox group.
  \item \textbf{filter-date-range} (lines~169--188): Native date inputs with
    range binding.
  \item \textbf{slider} (lines~190--203): \texttt{<sac-slider>} with
    min/max/step and format.
  \item \textbf{range-slider} (lines~205--231): Dual \texttt{<sac-slider>}
    instances for low/high bounds.
  \item \textbf{flex-container} (lines~233--249): \texttt{<sac-flex-container>}
    with recursive child widgets.
  \item \textbf{grid-container} (lines~251--265): \texttt{<sac-grid-container>}
    with column/row grid.
\end{enumerate}

\subsection{State Synchronization via STATE\_DELTA}

The \texttt{startStateSync()} method (lines~754--794) subscribes to an AG-UI
run with the sentinel message \texttt{\_\_state\_sync\_\_} and filters for
\texttt{STATE\_DELTA} and \texttt{CUSTOM} events. A generation counter
(\texttt{stateSyncGeneration}, line~418) prevents stale updates from
superseded subscriptions:

\begin{lstlisting}[language={[Sharp]C},caption={Stale-request guard (sac-ai-data-widget.component.ts:761--778)}]
const generation = ++this.stateSyncGeneration;

this.stateSub = this.agUiService
  .run({ message: '__state_sync__', ... })
  .pipe(
    filter((event): event is AgUiStateDeltaEvent
        | AgUiCustomEvent =>
      event.type === 'STATE_DELTA'
      || event.type === 'CUSTOM',
    ),
  )
  .subscribe({
    next: (event) => {
      if (generation !== this.stateSyncGeneration) {
        return; // discard stale events
      }
      ...
    },
  });
\end{lstlisting}

\subsection{Schema Normalization and Recursion Depth}

The \texttt{normalizeSchema()} method (lines~1137--1163) recursively normalizes
widget schemas. Children are only recursed when \texttt{depth < MAX\_CHILDREN\_DEPTH}
(line~1155), preventing unbounded nesting:

\begin{lstlisting}[language={[Sharp]C},caption={Depth-bounded normalization (sac-ai-data-widget.component.ts:1155--1157)}]
children: depth < MAX_CHILDREN_DEPTH
  ? schema.children?.map((child) =>
      this.normalizeSchema(child, depth + 1))
  : undefined,
\end{lstlisting}

\subsection{Chart Feeds and Table Columns}

Chart feed construction (lines~927--945) maps dimensions to
\texttt{Feed.CategoryAxis} and \texttt{Feed.Color}, and measures to
\texttt{Feed.ValueAxis}. Table column generation (lines~947--958) derives
column metadata from either a live \texttt{ResultSet} or the schema's dimension
and measure arrays. KPI value extraction (lines~1035--1048) reads the first
measure cell from the result set and computes trend direction by comparing
against the second row.

% =============================================================================
\section{Chat Panel with Human-in-the-Loop}
\label{sec:chat}
% =============================================================================

The \texttt{SacAiChatPanelComponent} (file:
\texttt{libs/sac-ai-widget/chat/sac-ai-chat-panel.component.ts}) provides the
conversational interface with streaming, tool confirmation, and audit logging.

\subsection{Streaming Chat}

The \texttt{send()} method (lines~618--665) creates a user message, pushes an
empty streaming assistant message, then subscribes to \texttt{agUiService.run()}.
The \texttt{handleEvent()} method (lines~775--821) processes events:

\begin{itemize}
  \item \texttt{TEXT\_MESSAGE\_CONTENT} appends delta text to the assistant
    message (line~779).
  \item \texttt{TOOL\_CALL\_START} creates a pending tool call buffer
    (line~786).
  \item \texttt{TOOL\_CALL\_ARGS} accumulates argument JSON fragments
    (line~800).
  \item \texttt{TOOL\_CALL\_END} triggers \texttt{executeTool()} (line~813).
\end{itemize}

\subsection{Tool Confirmation Queue}

The \texttt{queueConfirmation()} method (lines~883--902) manages a FIFO queue
of pending approvals. When a tool requires review (determined by
\texttt{toolDispatch.getConfirmationReview()}, line~837), the confirmation is
either set as \texttt{activeConfirmation} or pushed to
\texttt{queuedConfirmations}. The \texttt{advanceConfirmationQueue()} method
(lines~904--907) shifts the next pending review after approval or rejection.

\subsection{Approval and Rejection Workflow}

\begin{description}
  \item[\texttt{approveActiveConfirmation()}] (lines~689--737) Executes the
    reviewed tool via \texttt{toolDispatch.execute()}, dispatches the result
    back via \texttt{agUiService.dispatchToolResult()}, and records audit
    entries for both success and failure paths.
  \item[\texttt{rejectActiveConfirmation()}] (lines~739--761) Posts a rejection
    result with \texttt{code: 'USER\_REJECTED'} and advances the queue.
\end{description}

\subsection{Audit and Replay Logging}

Every significant action is double-logged through \texttt{recordAudit()} and
\texttt{recordReplay()} (lines~932--947), which delegate to
\texttt{SacAiSessionService}. The template renders both an audit timeline
(lines~178--203, ``Recent activity'') and a replay timeline (lines~156--176,
``Replay timeline'') with status badges and timestamps.

% =============================================================================
\section{Custom Element Entry Point}
\label{sec:entry}
% =============================================================================

The file \texttt{libs/sac-ai-widget/sac-ai-widget.entry.ts} registers the
\texttt{<sac-ai-widget>} custom element and manages the SAC Custom Widget
lifecycle.

\subsection{SAC Lifecycle Callbacks}

Three callbacks implement the SAC Custom Widget contract:

\begin{description}
  \item[\texttt{onCustomWidgetBeforeUpdate()}] (lines~73--75) Merges incoming
    designer properties (bearer token, backend URL, model ID, widget type)
    into a local \texttt{props} object.
  \item[\texttt{onCustomWidgetAfterUpdate()}] (lines~78--94) Triggers
    bootstrap on first call or when \texttt{capBackendUrl}/\texttt{tenantUrl}
    changes. Subsequent calls sync token and inputs without re-bootstrapping.
  \item[\texttt{onCustomWidgetDestroy()}] (lines~97--99) Destroys the Angular
    application reference.
\end{description}

\subsection{Angular Bootstrap in Shadow DOM}

The \texttt{bootstrap()} method (lines~105--166) creates an Angular application
via \texttt{createApplication()} with configured providers including
\texttt{SAC\_AI\_BACKEND\_URL}, \texttt{SAC\_TENANT\_URL}, and
\texttt{SAC\_MODEL\_ID} injection tokens.

The DOM structure is built inside a shadow root (line~133) for style isolation.
Two components are manually created via \texttt{createComponent()}:

\begin{lstlisting}[language={[Sharp]C},caption={Two-column layout (sac-ai-widget.entry.ts:140--161)}]
// Chat panel (left column, 320px)
const chatHost = document.createElement('div');
chatHost.style.cssText =
  'width:320px;min-width:280px;...flex-shrink:0;';
host.appendChild(chatHost);

// Data widget (right column, flex:1)
const dataHost = document.createElement('div');
dataHost.style.cssText = 'flex:1;overflow:hidden;';
host.appendChild(dataHost);

this.chatRef = createComponent(SacAiChatPanelComponent,
  { environmentInjector: injector, hostElement: chatHost });
this.dataRef = createComponent(SacAiDataWidgetComponent,
  { environmentInjector: injector, hostElement: dataHost });
\end{lstlisting}

\subsection{Custom Element Registration}

The element is registered at module load (lines~214--216):

\begin{lstlisting}[language={[Sharp]C},caption={Custom element registration (sac-ai-widget.entry.ts:214--216)}]
if (!customElements.get('sac-ai-widget')) {
  customElements.define('sac-ai-widget', SacAiWidgetEntry);
}
\end{lstlisting}

\noindent The guard prevents duplicate registration when the script is loaded
multiple times in an SAC story.

% =============================================================================
\section{Session and Trace Management}
\label{sec:session}
% =============================================================================

The \texttt{SacAiSessionService} (file:
\texttt{libs/sac-ai-widget/session/sac-ai-session.service.ts}) manages thread
identity, W3C trace context, and audit/replay logs.

\subsection{Thread ID Generation}

Thread IDs follow the pattern \texttt{sac-\{timestamp\}-\{random\}}
(line~118):

\begin{lstlisting}[language={[Sharp]C},caption={Thread ID format (sac-ai-session.service.ts:117--119)}]
private generateThreadId(): string {
  return `sac-${Date.now()}-${
    Math.random().toString(36).slice(2, 8)}`;
}
\end{lstlisting}

\subsection{W3C Traceparent (128-bit Hex)}

The \texttt{generateTraceId()} method (lines~121--126) produces a
W3C-compatible 32-character hexadecimal trace ID using the Web Crypto API:

\begin{lstlisting}[language={[Sharp]C},caption={Trace ID generation (sac-ai-session.service.ts:121--126)}]
private generateTraceId(): string {
  // W3C trace-context compatible: 32 hex chars (128-bit)
  const bytes = new Uint8Array(16);
  crypto.getRandomValues(bytes);
  return Array.from(bytes,
    b => b.toString(16).padStart(2, '0')).join('');
}
\end{lstlisting}

\subsection{Audit Entries}

The \texttt{SacAiAuditEntry} interface (lines~6--13) tracks event type, status
(\texttt{processing}, \texttt{approved}, \texttt{rejected}, \texttt{completed},
\texttt{error}), detail text, and trace ID. The \texttt{recordAudit()} method
(lines~67--83) prepends new entries and enforces a rolling limit of 20
(\texttt{auditLimit}, line~41).

\subsection{Replay Entries}

The \texttt{SacAiReplayEntry} interface (lines~15--32) adds a monotonic
\texttt{sequence} counter and a typed \texttt{kind} field with 10 possible
values:

\begin{lstlisting}[language={[Sharp]C},caption={Replay entry kinds (sac-ai-session.service.ts:19--31)}]
kind:
  | 'request.sent'
  | 'stream.chunk' | 'stream.complete' | 'stream.error'
  | 'tool.requested' | 'tool.result' | 'tool.error'
  | 'approval.required' | 'approval.queued'
  | 'approval.approved' | 'approval.rejected';
\end{lstlisting}

\noindent The replay log is capped at 100 entries (\texttt{replayLimit},
line~42). The \texttt{reset()} method (lines~59--65) clears both logs and
generates a fresh thread ID.

% =============================================================================
\section{Mangle Logic Rules}
\label{sec:mangle}
% =============================================================================

The Mangle file \texttt{mangle/sac\_widget.mg} encodes the widget taxonomy and
validation rules as Datalog-style predicates.

\subsection{Widget Taxonomy Predicates}

The 13 AI widget types are declared with a three-argument predicate
\texttt{sac\_ai\_widget\_type(type, category, role)} (lines~196--212):

\begin{lstlisting}[caption={Widget taxonomy facts (sac\_widget.mg:196--212)}]
sac_ai_widget_type("chart", "core", "visualization").
sac_ai_widget_type("table", "core", "visualization").
sac_ai_widget_type("kpi", "core", "visualization").

sac_ai_widget_type("filter-dropdown", "filter", "interactive").
sac_ai_widget_type("filter-checkbox", "filter", "interactive").
sac_ai_widget_type("filter-date-range", "filter", "interactive").

sac_ai_widget_type("slider", "slider", "interactive").
sac_ai_widget_type("range-slider", "slider", "interactive").

sac_ai_widget_type("text-block", "text", "display").
sac_ai_widget_type("heading", "text", "display").
sac_ai_widget_type("divider", "text", "display").

sac_ai_widget_type("grid-container", "layout", "container").
sac_ai_widget_type("flex-container", "layout", "container").
\end{lstlisting}

\subsection{Derived Rules and Type Guards}

Derived predicates encode structural constraints:

\begin{itemize}
  \item \texttt{sac\_can\_have\_children(WidgetType)} (lines~218--219) --
    only layout/container widgets may have children.
  \item \texttt{sac\_requires\_data\_binding(WidgetType)} (lines~222--223) --
    core visualization widgets require model binding.
  \item \texttt{sac\_requires\_dimension(WidgetType)} (lines~226--227) --
    filter widgets require a dimension.
  \item \texttt{sac\_is\_container\_widget}, \texttt{sac\_is\_filter\_widget},
    \texttt{sac\_is\_data\_widget} (lines~245--252) -- type guard predicates.
  \item \texttt{sac\_max\_children\_depth(8)} (line~215) -- mirrors
    \texttt{MAX\_CHILDREN\_DEPTH}.
\end{itemize}

\subsection{Filter Value Compatibility}

Four compatibility facts (lines~255--258) constrain which filter types accept
which value modes:

\begin{lstlisting}[caption={Filter compatibility rules (sac\_widget.mg:255--258)}]
sac_filter_value_compat("filter-dropdown", "SingleValue").
sac_filter_value_compat("filter-dropdown", "MultipleValue").
sac_filter_value_compat("filter-checkbox", "MultipleValue").
sac_filter_value_compat("filter-date-range", "RangeValue").
\end{lstlisting}

\subsection{API Endpoint Declarations}

REST endpoints for widgets, datasources, and planning are declared as
\texttt{api\_endpoint(domain, method, path)} facts (lines~265--282),
enabling static analysis of the API surface.

% =============================================================================
\section{Python Backend}
\label{sec:python}
% =============================================================================

The Python backend resides in \texttt{src/generativeUI/training-webcomponents-ngx/}.
Two files define the agent logic and the HTTP server.

\subsection{CircuitBreaker}

The \texttt{CircuitBreaker} class
(\texttt{agent/data\_cleaning\_agent.py}, lines~30--61) implements a three-state
circuit breaker (closed, open, half-open):

\begin{lstlisting}[language=Python,caption={CircuitBreaker (data\_cleaning\_agent.py:30--61)}]
class CircuitBreaker:
    def __init__(self, threshold: int = 5,
                 reset_timeout: float = 30.0):
        self.threshold = threshold
        self.reset_timeout = reset_timeout
        self._failure_count = 0
        self._last_failure_time: float = 0
        self._state = "closed"

    @property
    def state(self) -> str:
        if self._state == "open":
            if (time.monotonic()
                - self._last_failure_time
                    >= self.reset_timeout):
                self._state = "half-open"
        return self._state

    def allow_request(self) -> bool:
        s = self.state
        return s in ("closed", "half-open")
\end{lstlisting}

\noindent The breaker opens after 5 consecutive failures and automatically
transitions to half-open after 30 seconds.

\subsection{Exponential Backoff Retry}

The \texttt{retry\_with\_backoff()} function (lines~63--88) wraps an
async-compatible callable with exponential backoff:

\begin{lstlisting}[language=Python,caption={Exponential backoff (data\_cleaning\_agent.py:63--88)}]
async def retry_with_backoff(
    fn: Callable,
    max_retries: int = 3,
    base_delay: float = 0.5,
    circuit: Optional[CircuitBreaker] = None,
):
    last_error = None
    for attempt in range(max_retries):
        if circuit and not circuit.allow_request():
            break
        try:
            result = await fn()
            if circuit:
                circuit.record_success()
            return result
        except Exception as e:
            last_error = e
            if circuit:
                circuit.record_failure()
            if attempt < max_retries - 1:
                delay = base_delay * (2 ** attempt)
                await asyncio.sleep(delay)
    raise last_error if last_error else \
        RuntimeError("no attempts made")
\end{lstlisting}

\subsection{MangleQueryClient}

The \texttt{MangleQueryClient} (lines~91--141) sends JSON-RPC 2.0 requests to
the Mangle query service's MCP endpoint for field classification. It uses
\texttt{asyncio.to\_thread()} to run synchronous \texttt{urllib} calls in a
thread pool (line~134), and wraps each call with \texttt{retry\_with\_backoff()}
and a circuit breaker instance.

\subsection{ODataVocabularyDiscoveryClient}

The \texttt{ODataVocabularyDiscoveryClient} (lines~143--218) discovers field
schemas from the OData Vocabularies service. It exposes four methods:

\begin{enumerate}
  \item \texttt{discover\_entity\_schema()} (line~156) -- retrieves entity field
    definitions.
  \item \texttt{search\_vocabulary\_terms()} (line~160) -- searches vocabulary
    terms by query string.
  \item \texttt{get\_field\_classification()} (line~167) -- classifies a single
    field.
  \item \texttt{suggest\_annotations()} (lines~171--193) -- uses the OpenAI
    chat endpoint for annotation suggestions.
\end{enumerate}

\noindent All methods use circuit breakers and exponential backoff.

\subsection{FastAPI Server with Workflow SSE}

The HTTP server (\texttt{bin/api.py}) exposes a \texttt{/api/workflow/run}
endpoint (line~484) that returns a \texttt{StreamingResponse} with
\texttt{media\_type="text/event-stream"} (lines~755--763). The server includes
\texttt{TraceContextMiddleware} (lines~31--42) for W3C traceparent parsing.
Workflow events are emitted as SSE frames (line~271):

\begin{lstlisting}[language=Python,caption={SSE event emission (bin/api.py:270--271)}]
def _sse_event(payload: Dict[str, Any]) -> str:
    return f"data: {json.dumps(payload, default=str)}\n\n"
\end{lstlisting}

% =============================================================================
\section{Accessibility (WCAG AA Compliance)}
\label{sec:a11y}
% =============================================================================

The Generative UI system implements WCAG AA accessibility throughout all
interactive components.

\subsection{Semantic HTML}

The chat panel
(\texttt{sac-ai-chat-panel.component.ts}) uses semantic elements:

\begin{itemize}
  \item \texttt{<section aria-label="SAC AI Chat Assistant">} for the panel
    root (line~69).
  \item \texttt{<article>} elements for each chat message (line~84) with
    \texttt{aria-label} distinguishing user and assistant roles (line~88).
  \item \texttt{role="log"} with \texttt{aria-live="polite"} on the message
    container (lines~79--82).
  \item \texttt{role="form"} on the input row with labeling (line~206).
  \item \texttt{<label for="sac-chat-input">} with \texttt{sr-only} class
    (lines~207--209).
\end{itemize}

\noindent The heading component (\texttt{sac-text.component.ts}) renders
actual \texttt{<h1>} through \texttt{<h6>} elements via \texttt{ngSwitch}
(lines~39--45), preserving semantic heading hierarchy.

The divider component uses \texttt{<hr>} with \texttt{role="separator"}
(line~157) and \texttt{aria-label} (line~158).

\subsection{ARIA Attributes}

\begin{itemize}
  \item \textbf{Filter dropdown}: \texttt{aria-labelledby} linking to the
    label ID, \texttt{aria-describedby} for descriptions, and
    \texttt{aria-live="polite"} status announcements
    (\texttt{sac-filter.component.ts}, lines~55--70).
  \item \textbf{Filter checkbox}: \texttt{<fieldset>} with \texttt{<legend>}
    for proper grouping, \texttt{aria-describedby} for descriptions
    (\texttt{sac-filter.component.ts}, lines~150--168).
  \item \textbf{Slider}: Full \texttt{role="slider"} with
    \texttt{aria-valuemin}, \texttt{aria-valuemax}, \texttt{aria-valuenow},
    and \texttt{aria-valuetext} (\texttt{sac-slider.component.ts},
    lines~58--63).
  \item \textbf{Container widgets}: \texttt{role="group"} with
    \texttt{ariaLabel} on flex and grid containers
    (\texttt{sac-ai-data-widget.component.ts}, lines~240--241, 256--257).
  \item \textbf{Review panel}: \texttt{role="region"} with
    \texttt{aria-label="Planning action review"}
    (\texttt{sac-ai-chat-panel.component.ts}, lines~97--99).
    Error messages use \texttt{role="alert"} (line~134).
\end{itemize}

\subsection{SAP Fiori Design Tokens}

All color values use the SAP Fiori design token pattern
\texttt{var(--sapTokenName, fallback)}. No bare hex colors appear in
component styles. Examples from the codebase:

\begin{lstlisting}[caption={Fiori design token usage examples}]
/* Chat panel (sac-ai-chat-panel.component.ts:259-261) */
font-family: var(--sapFontFamily, '72', Arial, sans-serif);
font-size: var(--sapFontSize, 14px);
background: var(--sapBackgroundColor, #f7f7f7);

/* Data widget (sac-ai-data-widget.component.ts:294) */
font-family: var(--sapFontFamily, '72', Arial, sans-serif);
background: var(--sapField_Background, #fff);

/* Slider thumb (sac-slider.component.ts:93-94) */
background: var(--sapBrandColor, #0854a0);
border: 2px solid
  var(--sapContent_ContrastFocusColor, #fff);
\end{lstlisting}

\subsection{44px Touch Targets}

Interactive elements enforce a minimum 44px touch target:

\begin{itemize}
  \item Chat input: \texttt{min-height: 44px} (line~512 of
    \texttt{sac-ai-chat-panel.component.ts}).
  \item Send button: \texttt{min-height: 44px} (line~541).
  \item Review buttons: \texttt{min-height: 44px} (line~358).
  \item Date inputs: \texttt{min-height: 40px} with 8px padding
    (\texttt{sac-ai-data-widget.component.ts}, line~374).
\end{itemize}

\subsection{Keyboard Navigation and Reduced Motion}

The slider component (\texttt{sac-slider.component.ts}) implements enhanced
keyboard support (lines~154--179) with Home, End, PageUp, and PageDown keys
for coarse navigation. The chat panel and slider both respect
\texttt{@media (prefers-reduced-motion: reduce)} to disable animations
(chat panel line~464; slider line~113). A Windows High Contrast Mode media
query (\texttt{@media (forced-colors: active)}) is also present in the chat
panel (lines~572--578).

\subsection{Screen Reader Announcements}

Live regions (\texttt{aria-live="polite"}) with \texttt{sr-only} styling
provide non-visual feedback:

\begin{itemize}
  \item Chat panel: announces message sent, processing, response complete,
    and approval actions via the \texttt{announceMessage()} helper
    (lines~923--930).
  \item Filter dropdown: announces selection counts (line~70 of
    \texttt{sac-filter.component.ts}).
  \item Filter checkbox: announces select/deselect with total count
    (line~168 of \texttt{sac-filter.component.ts}).
  \item Slider: announces formatted value on change (line~72 of
    \texttt{sac-slider.component.ts}).
\end{itemize}

% =============================================================================
\section{Cross-Cutting Concerns}
\label{sec:crosscutting}
% =============================================================================

\subsection{Error Handling}

The system uses consistent error handling patterns:

\begin{enumerate}
  \item \textbf{SSE Stream}: Buffer overflow errors (line~198 of AG-UI
    service), HTTP errors (line~184), and abort handling (line~223).
  \item \textbf{Tool dispatch}: Every tool handler returns a \texttt{ToolResult}
    with \texttt{success: boolean} and optional \texttt{error: string}.
  \item \textbf{Python backend}: Three-tier fallback (Mangle rules,
    Elasticsearch cache, OData vocabulary) with circuit breakers on every
    external call.
  \item \textbf{Data widget}: Graceful degradation showing ``Live data
    unavailable'' status message on data binding failures (line~850).
\end{enumerate}

\subsection{Change Detection Strategy}

All Angular components use \texttt{ChangeDetectionStrategy.OnPush} to minimize
unnecessary change detection cycles. View updates are triggered explicitly via
\texttt{cdr.markForCheck()} after state mutations (e.g., line~633, 652, 781).
The data widget component includes a \texttt{flushView()} helper (lines~865--872)
that catches view-detached errors during destroy paths.

\subsection{Security Considerations}

\begin{itemize}
  \item Bearer tokens are stored in \texttt{SacAuthService} and transmitted
    via \texttt{Authorization} headers only over HTTPS.
  \item The Python backend's \texttt{TraceContextMiddleware} extracts trace
    context from request headers without executing any user-supplied code.
  \item The text block component sanitizes markdown HTML via Angular's
    \texttt{DomSanitizer} (\texttt{sac-text.component.ts}, line~138).
  \item CORS configuration in the Python API uses explicit origin allowlists
    and disables credentials when wildcards are present
    (\texttt{bin/api.py}, lines~48--58).
\end{itemize}

% =============================================================================
\section{Conclusion}
% =============================================================================

The SAP Analytics Cloud AI Widget system demonstrates a layered architecture
for generative UI: a type-safe 13-widget schema with recursive nesting, an
AG-UI protocol bridge with SSE streaming and W3C trace correlation, a tool
dispatch layer with human-in-the-loop approval for high-risk planning actions,
and a resilient Python backend with circuit breakers and multi-source field
classification. Accessibility is treated as a first-class concern throughout,
with semantic HTML, ARIA attributes, SAP Fiori design tokens, and 44px touch
targets ensuring WCAG AA compliance.

\vspace{1em}
\noindent\textbf{Source files referenced:}

\begin{longtable}{p{0.6\textwidth}p{0.35\textwidth}}
\toprule
\textbf{File} & \textbf{Section(s)} \\
\midrule
\texttt{libs/sac-ai-widget/types/sac-widget-schema.ts} & \ref{sec:schema} \\
\texttt{libs/sac-ai-widget/ag-ui/sac-ag-ui.service.ts} & \ref{sec:agui} \\
\texttt{libs/sac-ai-widget/chat/sac-tool-dispatch.service.ts} & \ref{sec:tools} \\
\texttt{libs/sac-ai-widget/data-widget/sac-ai-data-widget.component.ts} & \ref{sec:datawidget} \\
\texttt{libs/sac-ai-widget/chat/sac-ai-chat-panel.component.ts} & \ref{sec:chat}, \ref{sec:a11y} \\
\texttt{libs/sac-ai-widget/sac-ai-widget.entry.ts} & \ref{sec:entry} \\
\texttt{libs/sac-ai-widget/session/sac-ai-session.service.ts} & \ref{sec:session} \\
\texttt{mangle/sac\_widget.mg} & \ref{sec:mangle} \\
\texttt{src/generativeUI/training-webcomponents-ngx/packages/api-server/src/main.py} & \ref{sec:python} \\
\texttt{src/generativeUI/training-webcomponents-ngx/packages/api-server/src/train.py} & \ref{sec:python} \\
\texttt{libs/sac-ai-widget/components/sac-filter.component.ts} & \ref{sec:a11y} \\
\texttt{libs/sac-ai-widget/components/sac-slider.component.ts} & \ref{sec:a11y} \\
\texttt{libs/sac-ai-widget/components/sac-text.component.ts} & \ref{sec:a11y} \\
\bottomrule
\end{longtable}

\end{document}
